% =====================================================================
%  三、模型假设
% =====================================================================
\section{模型假设}
为便于建模求解，在合理范围内作如下假设：

\begin{enumerate}
    \item 示向度误差有界，$|\Delta\theta_i|\le\varepsilon$（本题 $\varepsilon
          =1^\circ$），且同一检测点的误差在观测期间固定不变，只有更换地点才呈现统计
          规律；故采用有界区间模型（最坏情形），所得区域保证包含真值。
    \item 检测点坐标与示向度读数精确已知；示向度以 $x$ 轴正向逆时针为正，取值于
          $[0^\circ,360^\circ)$。
    \item 干扰源位置在观测与清除期间固定不变，各源频道互不相同且不随时间变化。
    \item 测向机与干扰源位于同一平面，区域内无障碍遮挡。
    \item 全向源的覆盖角度为 $360^\circ$，定向源为定向方向两侧各 $90^\circ$（含边界），
          覆盖范围内信号均匀辐射、场强仅随距离衰减。
    \item 机器狗沿直线以 $5\,\mathrm{m/s}$ 匀速行进，忽略转向与加减速耗时；移动中
          不可检测，单次检测 $5\,\mathrm{s}$，任意两频道切换 $1\,\mathrm{s}$。
    \item 与干扰源距离不超过 $20\,\mathrm{m}$ 时，光学定位（$3\,\mathrm{s}$）与激光
          清除（$2\,\mathrm{s}$）必定成功且与覆盖角度无关；距离不超过 $5\,\mathrm{m}$
          且处于覆盖角度内时无法获得示向度，可直接定位并清除。
\end{enumerate}
